\section*{(4\textordfeminine\ aula) --- Produto de Operadores, Operador Adjunto, Hermitianos, Unitários e Autovalores (seções 1.6--1.8)}

\subsection*{Produto de Operadores em Forma Matricial}

Vimos que um operador $\hat{\Omega}$ corresponde à matriz com elementos $\Omega_{ij} = \bra{i}\hat{\Omega}\ket{j}$.

O \textbf{produto} $\hat{A}\hat{B}$ corresponde ao produto das matrizes:
\begin{align}
    (\hat{A}\hat{B})_{ij} = \bra{i}\hat{A}\hat{B}\ket{j}
    &= \bra{i}\hat{A}\underbrace{\left(\sum_k \ket{k}\bra{k}\right)}_{\hat{I}}\hat{B}\ket{j} \notag \\
    &= \sum_k \underbrace{\bra{i}\hat{A}\ket{k}}_{A_{ik}} \underbrace{\bra{k}\hat{B}\ket{j}}_{B_{kj}}
     = \sum_k A_{ik}\, B_{kj}
\end{align}
que é exatamente a regra do \textbf{produto matricial}. A correspondência operador $\leftrightarrow$ matriz preserva o produto.

\subsection*{Operador Adjunto (seção 1.7)}

\begin{shadequote}
    O \textbf{OPERADOR ADJUNTO} $\hat{\Omega}^\dagger$ de $\hat{\Omega}$ é definido por:
    \begin{equation}
        \bra{v}\hat{\Omega}^\dagger\ket{w} \equiv \bra{w}\hat{\Omega}\ket{v}^*
        \label{eq:adjunto_def}
    \end{equation}
    para todos $\ket{v}, \ket{w} \in \mathbb{V}^n$.
\end{shadequote}

O elemento de matriz $(ij)$ do adjunto é:
\begin{equation}
    (\hat{\Omega}^\dagger)_{ij} = \bra{i}\hat{\Omega}^\dagger\ket{j} = \bra{j}\hat{\Omega}\ket{i}^* = \Omega_{ji}^*
\end{equation}

A \textbf{matriz do adjunto é a transposta conjugada} da matriz original:
\begin{equation}
    \hat{\Omega}^\dagger \leftrightarrow (\Omega^T)^* = \Omega^\dagger
\end{equation}

Propriedades do adjunto:
\begin{align}
    (\hat{A} + \hat{B})^\dagger &= \hat{A}^\dagger + \hat{B}^\dagger \\
    (a\hat{A})^\dagger &= a^*\hat{A}^\dagger \\
    (\hat{A}\hat{B})^\dagger &= \hat{B}^\dagger\hat{A}^\dagger \label{eq:adjunto_produto} \\
    (\hat{A}^\dagger)^\dagger &= \hat{A}
\end{align}

\textbf{Prova de \eqref{eq:adjunto_produto}:}
Seja $\ket{z} = \hat{B}\ket{v}$; então, pela definição \eqref{eq:adjunto_def}:
\[
    \bra{v}(\hat{A}\hat{B})^\dagger\ket{w}
    = \bigl[\bra{w}\hat{A}\hat{B}\ket{v}\bigr]^*
    = \bigl[\bra{w}\hat{A}\ket{z}\bigr]^*
    = \bra{z}\hat{A}^\dagger\ket{w}
    = \bra{v}\hat{B}^\dagger\hat{A}^\dagger\ket{w} \quad \square
\]

\subsection*{Operadores Hermitianos e Unitários (seção 1.8)}

\begin{shadequote}
    \textbf{Operador Hermitiano} (ou auto-adjunto): $\hat{\Omega}^\dagger = \hat{\Omega}$

    Elementos de matriz: $\Omega_{ij} = \Omega_{ji}^*$ (a matriz é igual à sua transposta conjugada).
\end{shadequote}

\textbf{Exemplo:} $\displaystyle\Omega = \begin{bmatrix} 3 & 1+i \\ 1-i & 0 \end{bmatrix}$
é Hermitiana, pois $\Omega_{12} = 1+i = \Omega_{21}^*$.

\begin{shadequote}
    \textbf{Operador Unitário}: $\hat{U}^\dagger\hat{U} = \hat{U}\hat{U}^\dagger = \hat{I}$, ou equivalentemente $\hat{U}^{-1} = \hat{U}^\dagger$.
\end{shadequote}

\textbf{Exemplo:} $U = \dfrac{1}{\sqrt{2}}\begin{bmatrix} 1 & 1 \\ 1 & -1 \end{bmatrix}$ é unitária:
\[
    U^\dagger U = \frac{1}{2}\begin{bmatrix} 1 & 1 \\ 1 & -1 \end{bmatrix}\begin{bmatrix} 1 & 1 \\ 1 & -1 \end{bmatrix} = \begin{bmatrix} 1 & 0 \\ 0 & 1 \end{bmatrix} = I \quad \checkmark
\]

\begin{observacao}[]{}
    \textbf{Operador Anti-Hermitiano:} $\hat{\Omega}^\dagger = -\hat{\Omega}$.

    Qualquer operador pode ser decomposto na parte Hermitiana e Anti-Hermitiana:
    \[
        \hat{\Omega} = \underbrace{\frac{\hat{\Omega}+\hat{\Omega}^\dagger}{2}}_{\text{Hermitiana}} + \underbrace{\frac{\hat{\Omega}-\hat{\Omega}^\dagger}{2}}_{\text{Anti-Hermitiana}}
    \]
\end{observacao}

\subsection*{Autovalores e Autovetores (seção 1.8)}

\begin{shadequote}
    O \textbf{PROBLEMA DE AUTOVALORES} de um operador $\hat{\Omega}$: encontrar vetores $\ket{\omega}$ (não nulos) e números $\omega$ tais que
    \begin{equation}
        \hat{\Omega}\ket{\omega} = \omega\,\ket{\omega}
    \end{equation}
    Chamamos $\omega$ de \textbf{autovalor} e $\ket{\omega}$ de \textbf{autovetor} (ou autoestado) de $\hat{\Omega}$.
\end{shadequote}

Na MQ, grandezas físicas são representadas por operadores Hermitianos e os resultados possíveis de uma medição são os autovalores desses operadores.

Reescrevendo $(\hat{\Omega} - \omega\hat{I})\ket{\omega} = \ket{0}$: para existir solução não trivial, o operador deve ser singular. Em forma matricial:
\begin{equation}
    \det(\Omega - \omega\, I) = 0 \qquad \text{(equação característica)}
\end{equation}

Essa equação de grau $n$ em $\omega$ fornece os $n$ autovalores (com possíveis repetições). Para cada $\omega_k$, o autovetor é encontrado resolvendo o sistema linear $(\Omega - \omega_k I)\ket{\omega_k} = 0$.

\subsection*{Teoremas sobre Operadores Hermitianos}

\begin{shadequote}
    \textbf{Teorema 1:} Os autovalores de um operador Hermitiano são \textbf{REAIS}.
\end{shadequote}

\textit{Prova:} Seja $\hat{\Omega}^\dagger = \hat{\Omega}$ e $\hat{\Omega}\ket{\omega} = \omega\ket{\omega}$. Calcula-se $\bra{\omega}\hat{\Omega}\ket{\omega}$ de dois modos:
\[
    \bra{\omega}\hat{\Omega}\ket{\omega} = \omega\braket{\omega}{\omega}
    \quad\text{e}\quad
    \bra{\omega}\hat{\Omega}^\dagger\ket{\omega} = \bigl[\bra{\omega}\hat{\Omega}\ket{\omega}\bigr]^* = \omega^*\braket{\omega}{\omega}
\]
Como $\hat{\Omega} = \hat{\Omega}^\dagger$, as expressões são iguais. Como $\braket{\omega}{\omega} > 0$: $\omega = \omega^* \in \mathbb{R}$. $\square$

\begin{shadequote}
    \textbf{Teorema 2:} Autovetores de um operador Hermitiano associados a autovalores \textbf{DISTINTOS} são \textbf{ORTOGONAIS}.
\end{shadequote}

\textit{Prova:} Sejam $\omega_1 \neq \omega_2$ e $\hat{\Omega}\ket{\omega_i} = \omega_i\ket{\omega_i}$. Calculamos $\bra{\omega_1}\hat{\Omega}\ket{\omega_2}$:
\begin{align*}
    \bra{\omega_1}\hat{\Omega}\ket{\omega_2} &= \omega_2\braket{\omega_1}{\omega_2} \\
    \bra{\omega_1}\hat{\Omega}^\dagger\ket{\omega_2} &= \bigl[\bra{\omega_2}\hat{\Omega}\ket{\omega_1}\bigr]^* = \omega_1\braket{\omega_1}{\omega_2}
    \quad (\omega_1 \in \mathbb{R})
\end{align*}
Como $\hat{\Omega} = \hat{\Omega}^\dagger$: $(\omega_2 - \omega_1)\braket{\omega_1}{\omega_2} = 0$. Como $\omega_2 \neq \omega_1$: $\braket{\omega_1}{\omega_2} = 0$. $\square$

\subsection*{Autovalores Degenerados}

Dizemos que $\omega$ é \textbf{degenerado} com \textbf{degenerescência} $m$ se existem $m$ autovetores L.I.\ com esse autovalor. Eles formam um subespaço de dimensão $m$, o \textbf{autoespaço} de $\omega$.

Dentro do autoespaço pode-se escolher $m$ autovetores ortonormais, e qualquer combinação linear deles é autovetor de $\omega$:
\begin{equation}
    \hat{\Omega}\left(\sum_{a=1}^m c_a\ket{\omega,a}\right) = \omega\sum_{a=1}^m c_a\ket{\omega,a}
\end{equation}

\subsection*{Teoremas sobre Operadores Unitários}

\begin{shadequote}
    \textbf{Teorema 3:} Os autovalores de um operador Unitário têm módulo 1: $|\omega| = 1$, isto é, $\omega = e^{i\theta}$ para algum $\theta \in \mathbb{R}$.
\end{shadequote}

\textit{Prova:} De $\hat{U}^\dagger\hat{U} = \hat{I}$ e $\hat{U}\ket{\omega} = \omega\ket{\omega}$:
\[
    \braket{\omega}{\omega} = \bra{\omega}\hat{U}^\dagger\hat{U}\ket{\omega} = |\omega|^2\braket{\omega}{\omega}
    \quad \Longrightarrow \quad |\omega|^2 = 1 \quad \square
\]

\begin{shadequote}
    \textbf{Teorema 4:} Autovetores de um operador Unitário com autovalores DISTINTOS são ORTOGONAIS.
\end{shadequote}

A prova é análoga ao Teorema 2.

\subsection*{Exemplo: Operador Hermitiano com Autovalor Degenerado}

Considere, na base $\ket{1}, \ket{2}, \ket{3}$, o operador Hermitiano:
\begin{equation}
    \Omega = \begin{bmatrix} 2 & 1 & 0 \\ 1 & 2 & 0 \\ 0 & 0 & 3 \end{bmatrix}
\end{equation}

\textbf{Equação característica:}
\begin{equation}
    (3-\omega)\bigl[(2-\omega)^2 - 1\bigr] = 0
    \quad \Longrightarrow \quad \omega = 1,\quad \omega = 3\; (\text{degenerado, mult.\ 2})
\end{equation}

\textbf{Autovetor de $\omega = 1$:}
\[
    (\Omega - I)\ket{\omega_1} = 0 \quad \Longrightarrow \quad
    \begin{bmatrix} 1 & 1 & 0 \\ 1 & 1 & 0 \\ 0 & 0 & 2 \end{bmatrix}
    \begin{bmatrix} v_1 \\ v_2 \\ v_3 \end{bmatrix} = 0
    \quad \Longrightarrow \quad v_3 = 0,\; v_1 = -v_2
\]
\begin{equation}
    \ket{\omega=1} = \frac{1}{\sqrt{2}}\bigl(\ket{1} - \ket{2}\bigr)
\end{equation}

\textbf{Autoespaço de $\omega = 3$} (dimensão 2):
\[
    (\Omega - 3I)\ket{\omega_3} = 0 \quad \Longrightarrow \quad
    \begin{bmatrix} -1 & 1 & 0 \\ 1 & -1 & 0 \\ 0 & 0 & 0 \end{bmatrix}
    \begin{bmatrix} v_1 \\ v_2 \\ v_3 \end{bmatrix} = 0
    \quad \Longrightarrow \quad v_1 = v_2,\; v_3 \text{ livre}
\]
Escolhe-se dois autovetores ortonormais dentro do autoespaço:
\begin{equation}
    \ket{\omega=3, a=1} = \frac{1}{\sqrt{2}}\bigl(\ket{1}+\ket{2}\bigr), \qquad
    \ket{\omega=3, a=2} = \ket{3}
\end{equation}

A base ortonormal de autovetores é
$\left\{\dfrac{1}{\sqrt{2}}(\ket{1}-\ket{2}),\ \dfrac{1}{\sqrt{2}}(\ket{1}+\ket{2}),\ \ket{3}\right\}$,
e nessa base $\Omega$ é diagonal:
\begin{equation}
    \Omega \leftrightarrow \begin{bmatrix} 1 & 0 & 0 \\ 0 & 3 & 0 \\ 0 & 0 & 3 \end{bmatrix}
\end{equation}
